\section{The Split Likelihood Ratio Test and Conditional Independence Testing for Conditional Random Variables}


\begin{comment}
\subsubsection{The Split Likelihood Ratio Test}

In the following we outline the general framework of the \emph{split likelihood ratio test} from \cite{WRB20}.

\begin{Prn}[Split likelihood ratio test, see \cite{WRB20,SD22}]
    Let $\Xbf=\lC x^{(1)},\dots,x^{(N)} \rC \ins \Xcal$ be i.i.d.\ data from a probability distribution $P \in \Pcal(\Xcal)$.
 Let $\Hcal_0 = \lC P_\theta \st \theta \in \Theta_0 \rC $ and 
 $\Hcal_1 = \lC P_\theta \st \theta \in \Theta_1 \rC \ins \Pcal(\Xcal)$ be two disjoint hypothesis classes given by
 densities $x \mapsto p(x|\theta)$ w.r.t.\ a reference measure $\mu$ on $\Xcal$.
 We want to test if $P \in \Hcal_0$ vs.\ $P \in \Hcal_1$.
 \begin{enumerate}
     \item Split the data roughly 1:1 into two disjoint sets: $\Xbf = \Xbf_0 \dcup \Xbf_1$ (split unevenly, e.g.\ 1:2, if $\dim \Theta_0 \ll \dim \Theta_1$).
     \item Train a model $x \mapsto \hat p(x|\Xbf_1)$ using any method on $\Xbf_1$ and $\Hcal_1$, e.g.\ via maximum likelihood fitting: $\hat p(x|\Xbf_1) = p(x|\hat \theta_1(\Xbf_1))$:
         \[ \hat \theta_1(\Xbf_1) := \argmax_{\theta \in \Theta_1}  p(\Xbf_1|\theta), \]
         where we make the abbreviation:
         \[ p(\Xbf_i|\theta) := \prod_{x \in \Xbf_i} p(x|\theta). \]
     \item Train a model $x \mapsto p(x|\hat \theta_0(\Xbf_0))$ via maximum likelihood fitting on $\Xbf_0$ and $\Hcal_0$ :
         \[ \hat \theta_0(\Xbf_0) := \argmax_{\theta \in \Theta_0}  p(\Xbf_0|\theta). \]
         Note that we have: 
         \[  p(\Xbf_0|\hat \theta_0(\Xbf_0)) = \sup_{\theta \in \Theta_0} p(\Xbf_0|\theta).   \]
     \item Consider the test statistic:
         \[ E := \frac{ \hat p(\Xbf_0|\Xbf_1)}{p(\Xbf_0|\hat \theta_0(\Xbf_0))} = \frac{ \hat p(\Xbf_0|\Xbf_1)}{\sup_{\theta \in \Theta_0} p(\Xbf_0|\theta)},  \]
            where both densities are evaluated on $\Xbf_0$ by taking the product of the densities for each point.
        \item      Further note that $E$ is an \emph{E-variable} for $\Hcal_0$, i.e.:
     \begin{align*} \sup_{\theta \in \Theta_0} \E_\theta\lB E \rB  \le 1.
     \end{align*}
       Indeed, for every $\theta \in \Theta_0$, i.e.\ assuming $\Xbf \sim P_\theta \in \Hcal_0$ i.i.d., we have:
         \begin{align*} \E_\theta\lB E \rB 
         &= \int\int \frac{ \hat p(x_0|x_1)}{p(x_0|\hat \theta_0(x_0))} \, p(x_0,x_1|\theta) \, d\mu(x_0)\, d\mu(x_1) \\
         &\stackrel{\hat \theta_0 \text{ MLE}}{\le} \int\int \frac{ \hat p(x_0|x_1)}{p(x_0|\theta)} \, p(x_0,x_1|\theta) \, d\mu(x_0)\, d\mu(x_1) \\ 
         &\stackrel{\Xbf_0 \Indep \Xbf_1|\theta}{=} \int\int \frac{ \hat p(x_0|x_1)}{p(x_0|\theta)} \, p(x_0|\theta) \, p(x_1|\theta) \, d\mu(x_0)\, d\mu(x_1) \\
         & = \int \lp \int \hat p(x_0|x_1) \, d\mu(x_0)  \rp \, p(x_1|\theta)  \, d\mu(x_1) \\
         & = 1.
     \end{align*}
 \item Pick a significance level $\alpha \in [0,1]$\footnote{Historically: $\alpha=0.05$, 
     in physics: $\alpha = \frac{1}{3488556}= 3 \cdot 10^{-7}$ (5$\sigma$-rule).} and reject $\Hcal_0$ in favor of $\Hcal_1$ if:
         \[ \Lambda:=\log E \ge - \log \alpha. \]
         We call $-\log_{2} \alpha$ the \emph{bits of safety} (against type-I-error) 
         and $\Lambda/\log 2$ the \emph{evidential bits of safety} (for rejecting $\Hcal_0$ in favor of $\Hcal_1$).
         If we use $\log_e$ we would say \emph{nats of safety}. For $\log_{10}$ it was proposed to use the term \emph{nines of safety}\footnote{
         \texttt{https://terrytao.wordpress.com/2021/10/03/nines-of-safety-a-proposed-unit-of-measurement-of-risk}}.
     \item The type-I-error is then bounded by $\alpha$. 
     \item[]  Indeed, for every $\alpha \in [0,1]$ and $\theta \in \Theta_0$ we have with help of the Markov inequality:
         \begin{align*} P_\theta\lp \Lambda \ge - \log \alpha  \rp &= P_\theta \lp E \ge \alpha^{-1} \rp 
             \le \frac{\E_\theta\lB E \rB }{\alpha^{-1}} = \E_\theta\lB E \rB \cdot \alpha \stackrel{\E_\theta\lB E \rB  \le 1}{\le} \alpha.
         \end{align*}
 \end{enumerate}
\end{Prn}

\subsubsection{Conditional Independence Testing for Random Variables}

Now again consider three categorical random variables $X,Y,Z$ taking values in spaces $\Xcal$, $\Ycal$
and $\Zcal$, resp., with $2 \le K:= |\Xcal| < \infty$, $2 \le L:=|\Ycal| < \infty$ and $1 \le M:=|\Zcal| < \infty$, 
and joint distribution $P(X,Y,Z)$.

Let $\Dbf=\lC (x_n,y_n,z_n) \st n =1, \dots, N \rC$ be a data set of size $N$ drawn i.i.d.\ from $P(X,Y,Z)$.
We want to test:
\[ H_0:\, X \Indep_{P(X,Y,Z)} Y \given Z, \qquad \text{v.s.} \qquad H_1:\,X \nIndep_{P(X,Y,Z)} Y \given Z. \]
We will parameterize the mass function of $P(X,Y,Z|\theta)$ as:
\[ p(x,y,z|\theta) = \theta_{x,y,z}, \qquad \theta = \lp \theta_{x,y,z} \rp_{(x,y,z) \in \Xcal\times\Ycal\times\Zcal}, \]
for $\theta \in \Theta^{(0)}$, $\theta \in \Theta^{(1)}$, resp.

Split the data between 1:1 to 1:2: $\Dbf = \Dbf^{(0)} \dcup \Dbf^{(1)}$. 
We then put, similar to before, $i=0,1$:
\begin{align*}
    N^{(i)} & := \# \Dbf^{(i)},\\
    N^{(i)}_{x,y,z} &:=  \sum_{(x',y',z') \in \Dbf^{(i)}} \I_{(x,y,z)}(x',y',z'),\\
    N^{(i)}_{x,+,z} &:=  \sum_{(x',y',z') \in \Dbf^{(i)}} \I_{(x,z)}(x',z') &&= \sum_{y \in \Ycal} N^{(i)}_{x,y,z} ,\\
    N^{(i)}_{+,y,z} &:=  \sum_{(x',y',z') \in \Dbf^{(i)}} \I_{(y,z)}(y',z') &&= \sum_{x \in \Xcal} N^{(i)}_{x,y,z} ,\\
    N^{(i)}_{+,+,z} &:=  \sum_{(x',y',z') \in \Dbf^{(i)}} \I_{z}(z') &&= \sum_{x \in \Xcal} N^{(i)}_{x,+,z} ,\\
    N^{(i)}_{x,+,+} &:=  \sum_{(x',y',z') \in \Dbf^{(i)}} \I_{x}(x') &&= \sum_{z \in \Zcal} N^{(i)}_{x,+,z}.
\end{align*}
The maximum likelihood objective then leads to the following frequency estimators:
\[ \hat \theta^{(1)}_{x,y,z} := \frac{N^{(1)}_{x,y,z}}{N^{(1)}}, \qquad \hat \theta^{(0)}_{x,y,z} := \hat \theta^{(0)}_{x|z} \cdot \hat \theta^{(0)}_{y,z}, \quad \text{ with: } \quad
\hat \theta^{(0)}_{y,z} := \frac{N^{(0)}_{+,y,z}}{N^{(0)}}, \qquad \hat \theta^{(0)}_{x|z} := \frac{N^{(0)}_{x,+,z}}{N^{(0)}_{+,+,z}},  \]
or, more precisely, to avoid division by zero:
\[
\hat \theta^{(0)}_{x|z} :=  \left\{ 
    \begin{array}{lll}
     \frac{N^{(0)}_{x,+,z}}{N^{(0)}_{+,+,z}}  & \text{ if } & N^{(0)}_{+,+,z} \neq 0,\\
     \frac{N^{(0)}_{x,+,+}}{N^{(0)}} \cdot  & \text{ if } & N^{(0)}_{+,+,z} = 0.
\end{array} \right.
\]
Note that in the split likelihood ratio test we have some freedom to fit $\hat \theta^{(1)}$. To avoid (or enforce) the vanishing of some
$\hat \theta^{(1)}_{x,y,z}$, we could also use a regularized version by adding (or subtracting, resp.) small ``pseudo-counts'' $\epsilon \in \R$:
\begin{align*}
    N^{(1),\epsilon}_{x,y,z} := \max(0,N^{(1)}_{x,y,z} + \epsilon),\qquad N^{(1),\epsilon} := \sum_{(x,y,z)\in\Xcal\times\Ycal\times\Zcal}N^{(1),\epsilon}_{x,y,z},\qquad
\hat \theta^{(1)}_{x,y,z} :=  \frac{N^{(1),\epsilon}_{x,y,z}}{N^{(1),\epsilon}}.
\end{align*}
We now put for $i=0,1$:
\[ \hat \theta^{(i)} := \lp \hat \theta^{(i)}_{x,y,z} \rp_{(x,y,z)  \in \Xcal\times \Ycal\times \Zcal}.  \]
We can then write the estimated mass functions as:
\[ p(x,y,z|\hat \theta^{(i)}) = \prod_{(x',y',z')\in\Xcal\times\Ycal\times\Zcal} \lp\hat \theta^{(i)}_{x',y',z'}\rp^{\I_{(x',y',z')}(x,y,z)}.\]
Plugging $\Dbf^{(j)}$ in we get:
\begin{align*}
    p(\Dbf^{(j)}|\hat \theta^{(i)}) 
   & = \prod_{(x,y,z) \in \Dbf^{(j)}} \prod_{(x',y',z')\in\Xcal\times\Ycal\times\Zcal} \lp\hat \theta^{(i)}_{x',y',z'}\rp^{\I_{(x',y',z')}(x,y,z)}\\
   & =  \prod_{(x',y',z')\in\Xcal\times\Ycal\times\Zcal} \lp\hat \theta^{(i)}_{x',y',z'}\rp^{N_{x',y',z'}^{(j)}}.
\end{align*}
So the test statistic for the split likelihood ratio test is:
\begin{align*} 
    E &= \prod_{(x,y,z)\in\Xcal\times\Ycal\times\Zcal} \lp  \frac{\hat \theta^{(1)}_{x,y,z}}{\hat \theta^{(0)}_{x,y,z}} \rp^{N_{x,y,z}^{(0)}}.
\end{align*}
and thus:
\begin{align*} 
    \Lambda = \log E &= \sum_{(x,y,z)\in\Xcal\times\Ycal\times\Zcal} N_{x,y,z}^{(0)} \cdot \log \frac{\hat \theta^{(1)}_{x,y,z}}{\hat \theta^{(0)}_{x,y,z}} \\
    &=  \sum_{(x,y,z)\in\Xcal\times\Ycal\times\Zcal} N_{x,y,z}^{(0)} \cdot \log \frac{N^{(1),\epsilon}_{x,y,z} \cdot N^{(0)}_{+,+,z} \cdot N^{(0)}}
    {N^{(1),\epsilon} \cdot N^{(0)}_{x,+,z} \cdot N^{(0)}_{+,y,z}}. 
\end{align*}
For a significance level $\alpha \in [0,1]$ we declare $X$ dependent on $Y$ conditioned on $Z$ if:
\[ \Lambda \ge -\log \alpha. \]
As stated before, the type-I-error will be bounded by $\alpha$.

\begin{Rem}
    On the one hand, the split likelihood ratio test is usually less powerful than the classical likelihood ratio test, due to two reasons:
    \begin{enumerate}
        \item If the split is given by the ratio $\rho := \frac{\#\Dbf_0}{\#\Dbf} < 1$ then:
            \[ N_{x,y,z}^{(0)}  \approx \rho \cdot N_{x,y,z}.   \]
            So $\Lambda$ is roughly by a factor of $\rho$ smaller
            than $\frac{1}{2} \cdot G_N$ (note the correction for the conventional factor of $2$ inside $G_N$), 
            where the latter is fitted on the whole dataset $\Dbf$.
        \item The corresponding thresholds are typically related like:
    \[ \frac{1}{2} \cdot \chi^2_{\nu,1-\alpha} < - \log \alpha. \]
    \end{enumerate}
    So if $H_1$ was true, then $\frac{1}{2} G_N$ would more quickly surpass the threshold of $\frac{1}{2} \cdot \chi^2_{\nu,1-\alpha}$ than $\Lambda$ would surpass $- \log \alpha$.

    On the other hand, the advantage of the split likelihood ratio test is that the test statistic is based on an \emph{E-variable}.
    It thus inherits all the perks from safe testing, see \cite{GHK20}, like optional stopping and optional continuation.    
    Furthermore, the type-I-error control holds for all (finite) sample sizes and not just asymptotically.
    \PF{maybe cite other papers}
\end{Rem}

\end{comment}


\subsubsection{The Split Likelihood Ratio Test for Conditional Random Variables}

In the following we generalize the general framework of the \emph{split likelihood ratio test} from \cite{WRB20} to Markov kernels and conditional random variables. For this we first need to generalize the notion of E-variable, see \cite{GHK20}, to this scenario as well.

\begin{Def}[E-variable]
    \label{def:e-variable}
    Let $\Xcal$ and $\Tcal$ be two measurable spaces and $\Hcal_0=\lC P_\theta(X|T) \st \theta \in \Theta_0 \rC$
    be a class of Markov kernels: $\Tcal \dshto \Xcal$.
    A measurable map $E:\, \Xcal \times \Tcal \to \R$ is called \emph{E-variable for $\Hcal_0$} if $E \ge 0$ and 
    for all $\theta \in \Theta_0$ and all $t \in \Tcal$ we have:
    \[ \E_\theta[E|T=t] :=  \int E(x,t) \, P_\theta(X \in dx|T=t) \le 1.   \]
    %\[  \sup_{\theta \in \Theta_0} \E_\theta[E|T] \le 1.     \]
\end{Def}

\begin{Rem}
    Let $E$ be an E-variable for $\Hcal_0$. 
    Then we have for every $\theta \in \Theta_0$ and every probability distribution $Q(T) \in \Prob(\Tcal)$:
    \[ \E_{Q,\theta}[E] :=  \int E(x,t) \, P_\theta(X \in dx|T=t) \, Q(T \in dt) \le 1.   \]
    In other words, $E$ is an E-variable for:
    \[\tilde \Hcal_0 := \lC P_\theta(X|T)\otimes Q(T) \st \theta \in \Theta_0, Q(T) \in \Prob(\Tcal) \rC.\]
\end{Rem}

The main advantage of building test statistics on E-variables is that we have a non-asymptotic error-I-control:

\begin{Lem}
    \label{lem:e-variable-I-error}
    Let  $E:\, \Xcal \times \Tcal \to \R_{\ge0}$ be an E-variable for $\Hcal_0=\lC P_\theta(X|T) \st \theta \in \Theta_0 \rC$.
    Then for every $\theta \in \Theta_0$, $t \in \Tcal$ and $\alpha \in [0,1]$ we have:
    \[  P_\theta(E \ge \alpha^{-1}|T=t) \le \alpha,   \]
    where we put $0^{-1}:=\infty$.
    \begin{proof}
        This follows directly from the Markov inequality:
        \begin{align*}
            P_\theta(E \ge \alpha^{-1}|T=t) \le \frac{\E_\theta[E|T=t]}{\alpha^{-1}} = \alpha \cdot \E_\theta[E|T=t] 
            \stackrel{\E_\theta[E|T=t] \le 1 }{\le} \alpha.
        \end{align*}
        This shows the claim.
    \end{proof}
\end{Lem}


\begin{Prn}[Split likelihood ratio test for Markov kernels] $\,$
    \begin{enumerate}
\item Let $\Xcal$ and $\Tcal$ be two measurable spaces and $\Hcal_0=\lC P_\theta(X|T) \st \theta \in \Theta_0 \rC$
     and $\Hcal_1=\lC P_\theta(X|T) \st \theta \in \Theta_1 \rC$ 
    be two classes of Markov kernels: $\Tcal \dshto \Xcal$.
\item We assume that every $P_\theta(X|T)$ has a density w.r.t.\ a reference measure $\mu$ on $\Xcal$. 
      We denote the density in arguments as $p(x|t,\theta)$. % or $p_\theta(x|t)$.
\item Let $\Dbf = \lC (X_1,T_1),\dots,(X_N,T_N) \rC$ be a data set of i.i.d.\ variables from a Markov kernel $P(X|T)$, i.e.\
    we consider $\Xbf:=(X_1,\dots,X_N)$  to be sampled with input $\Tbf=(T_1,\dots,T_N)$ from the product Markov kernel:
    \[ P^{\otimes N}(\Xbf|\Tbf) := \bigotimes_{n=1}^N P(X_n|T_n), \]
    where each $P(X_n|T_n)$ is a copy of $P(X|T)$. 
\item Based on $\Dbf$ we want to test:
    \[ H_0:\; P(X|T) \in \Hcal_0 \qquad \text{ vs. } \qquad H_1:\; P(X|T) \in \Hcal_1. \]
\item Split the data into two disjoint sets $\Dbf = \Dbf^{(0)} \dcup \Dbf^{(1)}$ with corresponding $\Xbf^{(0)}$, $\Xbf^{(1)}$ 
    and $\Tbf^{(0)}$, $\Tbf^{(1)}$, resp. %Let $I_i:=\lC n \st (X_n,T_n) \in \Dbf^{(i)}  \rC$ the index set corresponding to $\Dbf^{(i)}$, $i=0,1$.
\item Take the maximum likelihood estimators for $i=0,1$, each on their own data set $\Dbf^{(i)}$ and parameter space $\Theta_i$:
    \[ \hat \theta^{(i)} := \hat \theta^{(i)}(\Dbf^{(i)}) := \argmax_{\theta \in \Theta_i}  p(\Xbf^{(i)}|\Tbf^{(i)},\theta) 
    = \argmax_{\theta \in \Theta_i} \prod_{(X_n,T_n) \in \Dbf^{(i)}} p(X_n|T_n,\theta). \]
\item Define the test statistic:
    \[ E := E(\Dbf) := \frac{p(\Xbf^{(0)}|\Tbf^{(0)},\hat \theta^{(1)}(\Dbf^{(1)}))}{p(\Xbf^{(0)}|\Tbf^{(0)},\hat \theta^{(0)}(\Dbf^{(0)}))} 
    = \prod_{(X_n,T_n) \in \Dbf^{(0)}} \frac{p(X_n|T_n,\hat\theta^{(1)})}{p(X_n|T_n,\hat\theta^{(0)})},  \]
    where both densities are evaluated on $\Dbf^{(0)}$.
\item Then $E$ is an E-variable for $\Hcal_0^{\otimes N} := \lC P_\theta^{\otimes N}(\Xbf|\Tbf) \st \theta \in \Theta_0 \rC$.
\item[] Indeed, we have for every $\theta \in \Theta_0$ and $\tbf \in \Tcal^N$:
    \end{enumerate}
    \begin{align*}
        &\quad \E_\theta^{\otimes N}[E|\Tbf=\tbf] \\ 
        & = \int\int \frac{p(\xbf^{(0)}|\tbf^{(0)},\hat \theta^{(1)}(\xbf^{(1)},\tbf^{(1)}))}{p(\xbf^{(0)}|\tbf^{(0)},\hat \theta^{(0)}(\xbf^{(0)},\tbf^{(0)}))} \cdot p(\xbf^{(0)}|\tbf^{(0)},\theta)\cdot p(\xbf^{(1)}|\tbf^{(1)},\theta) \,\mu(d\xbf^{(0)})\,\mu(d\xbf^{(1)}) \\ 
        & \stackrel{\hat \theta^{(0)} \text{ MLE}}{\le} \int\int \frac{p(\xbf^{(0)}|\tbf^{(0)},\hat \theta^{(1)}(\xbf^{(1)},\tbf^{(1)}))}{p(\xbf^{(0)}|\tbf^{(0)},\theta)} \cdot p(\xbf^{(0)}|\tbf^{(0)},\theta)\cdot p(\xbf^{(1)}|\tbf^{(1)},\theta) \,\mu(d\xbf^{(0)})\,\mu(d\xbf^{(1)}) \\ 
        & =  \int \lp\int p(\xbf^{(0)}|\tbf^{(0)},\hat \theta^{(1)}(\xbf^{(1)},\tbf^{(1)})) \,\mu(d\xbf^{(0)}) \rp p(\xbf^{(1)}|\tbf^{(1)},\theta) \,\mu(d\xbf^{(1)}) \\
        & =1.
    \end{align*}
    \begin{enumerate}[resume]
\item Pick a significance level $\alpha \in [0,1]$ and reject $\Hcal_0$ in favor of $\Hcal_1$ if:
         \[ \Lambda:=\log E \ge - \log \alpha. \] 
 \item The type-I-error is then bounded by $\alpha$ via Lemma \ref{lem:e-variable-I-error} 
\item[]  Indeed, for every $\alpha \in [0,1]$, $\theta \in \Theta_0$, $\tbf \in \Tcal^N$ we have: 
        %with help of the Markov inequality:
         \begin{align*} 
             P_\theta^{\otimes N}\lp \Lambda \ge - \log \alpha |\Tbf=\tbf  \rp 
             &= P_\theta^{\otimes N} \lp E \ge \alpha^{-1} |\Tbf=\tbf \rp \le \alpha. %\\ 
             %&\le \frac{\E_\theta^{\otimes N} \lB E|\Tbf=\tbf \rB }{\alpha^{-1}} \\
             %& \le \alpha.
         \end{align*}
\item[] Note that this implies that for every \emph{joint} distribution
    $Q(\Tbf) \in \Prob(\Tcal^N)$ and $\theta \in \Theta_0$ we have the following type-I-error control:
    \[ Q_\theta \lp \Lambda \ge - \log \alpha  \rp \le \alpha, \]
    where:
    \[ Q_\theta(\Xbf,\Tbf) := P_\theta^{\otimes N}(\Xbf|\Tbf) \otimes Q(\Tbf).\]
    \end{enumerate}
\end{Prn}

\begin{Rem}[Power/type-II-error]
    Note that while the maximum likelihood estimation as the fitting choice for $\hat \theta^{(0)}(\Dbf^{(0)})$
    is somewhat necessary for the arguments about the type-I-error control to work,
    we could use other fitting procedures for $\hat \theta^{(1)}(\Dbf^{(1)})$. 
    While this would effect the power/type-II-error control of the test, 
    the arguments above would still work to achieve the same type-I-error control.
    To achieve a high power for the split likelihood ratio test the rule of thumb is that we require $p(X|T,\hat \theta^{(1)})$ 
    to take on big values whenever $(X,T) \sim P_\theta(X|T) \in \Hcal_1$.
    Since we only see data points from the ``true'' underlying Markov kernel $P(X|T)$ the maximum likelihood 
    fit for $\hat \theta^{(1)}(\Dbf^{(1)})$ on $\Hcal_1$ and $\Dbf^{(1)} \sim P(X|T)$ generally is a reasonable choice.
    Some regularization might be beneficial to make sure the fitted model $P_{\hat \theta^{(1)}}(X|T)$ generalizes well to 
    the unseen data set $\Dbf^{(0)} \sim P(X|T)$.

\begin{comment}
    If we stick to the maximum likelihood estimators we can also define the variable:
    \[ \Lambda_1 := \Lambda_1(\Dbf) := \log \frac{p(\Xbf^{(1)}|\Tbf^{(1)},\hat \theta^{(1)}(\Dbf^{(1)}))}{p(\Xbf^{(1)}|\Tbf^{(1)},\hat \theta^{(0)}(\Dbf^{(0)}))}. \]
    By a symmetric argument to the above we get for every power level $\beta \in [0,1]$, $\tbf \in \Tcal^N$,
    $\theta \in \Theta_1$ the following minimal power:
    %\[ P_\theta^{\otimes N}\lp -\Lambda_1 \ge  -\log (1-\beta) |\Tbf=\tbf  \rp \le 1-\beta. \]
    \[ P_\theta^{\otimes N}\lp \Lambda_1 >  \log (1-\beta) |\Tbf=\tbf  \rp \ge \beta.  \]
 In other words, the type-II-error for the test 
 ``reject $H_0$ in favor of $H_1$ if $\Lambda_1 > \log(1-\beta)$''
 is bounded by $1-\beta$.
 This means we could use a combination of $\Lambda_0:=\Lambda$ and $\Lambda_1$ as test statistic 
 to trade off between type-I-error control and power/type-II-error control.
\end{comment}
\end{Rem}


\subsubsection{Categorical Conditional Random Variables}

We are now going to use the tools from the last subsection to build a split likelihood ratio test for conditional independence testing for finite discrete conditional random variables and Markov kernels.

Now consider three categorical conditional random variables $X,Y,Z$ taking values in spaces $\Xcal$, $\Ycal$
and $\Zcal$, resp., with $2 \le  |\Xcal| < \infty$, $2 \le |\Ycal| < \infty$ and $1 \le |\Zcal| < \infty$, 
and input space $\Tcal$ with $1 \le |\Tcal| < \infty$, 
and joint Markov kernel  $P(X,Y,Z|T)$.

Let $\Dbf=\lC (X_n,Y_n,Z_n,T_n) \st n =1, \dots, N \rC$ be a data set of size $N$ drawn i.i.d.\ from $P(X,Y,Z|T)$.
Note that for each data point $(X_n,Y_n,Z_n)$ we also recorded the input $T_n$ that generated that point.
We want to test:
\[ H_0:\, X \Indep_{P(X,Y,Z|T)} Y \given Z, \qquad \text{vs.} \qquad H_1:\,X \nIndep_{P(X,Y,Z|T)} Y \given Z. \]
We will parameterize the mass function of $P_\theta(X,Y,Z|T)$ as:
\[ p(x,y,z|t,\theta) = \theta_{x,y,z|t}, \qquad \theta = \lp \theta_{x,y,z|t} \rp_{(x,y,z,t) \in \Xcal\times\Ycal\times\Zcal\times\Tcal}, \]
for $\theta \in \Theta^{(0)}$ and $\theta \in \Theta^{(1)}$, resp., where additional normalization and factorization constraints are imposed. 
More clearly, we use the following parameter spaces:
\begin{align*}
    \Theta^{(0)}_{X|Z} &:= \lC \theta_{X|Z} =\lp \theta_{x|z} \rp_{(x,z) \in \Xcal\times\Zcal}  \in \prod_{(x,z) \in \Xcal\times\Zcal} [0,1] \st \forall z \in \Zcal. \, \sum_{x \in \Xcal} \theta_{x|z} =1   \rC,\\
    \Theta^{(0)}_{Y,Z|T} &:= \lC \theta_{Y,Z|T} =\lp \theta_{y,z|t} \rp_{(y,z,t) \in \Ycal\times\Zcal\times\Tcal}  \in \prod_{(y,z,t) \in \Ycal\times\Zcal\times\Tcal} [0,1] \st \forall t \in \Tcal. \, \sum_{(y,z) \in \Ycal\times\Zcal} \theta_{y,z|t} =1   \rC,\\
   % \Theta^{(0)} &:= \Theta^{(0)}_{X|Z} \times \Theta^{(0)}_{Y,Z|T},\\
   % \Theta^{(0)} &:= \lC \theta = \lp \theta_{x|z} \cdot \theta_{y,z|t} \rp_{(x,y,z,t) \in\Wcal} \st \theta_{X|Z} \in \Theta^{(0)}_{X|Z},\theta_{Y,Z|T} \in \Theta^{(0)}_{Y,Z|T}  \rC,\\
    \Theta^{(0)} &:= \lC \theta = \lp \theta_{x|z} \cdot \theta_{y,z|t} \rp_{(x,y,z,t) \in\Wcal} \st \lp \theta_{x|z} \rp_{(x,z) \in \Xcal\times\Zcal}  \in \Theta^{(0)}_{X|Z},\lp \theta_{y,z|t} \rp_{(y,z,t) \in \Ycal\times\Zcal\times\Tcal} \in \Theta^{(0)}_{Y,Z|T}  \rC,\\
    \Theta^{(1)} &:= \lC \theta = \lp \theta_{x,y,z|t} \rp_{(x,y,z,t) \in\Wcal} \in \prod_{(x,y,z,t) \in\Wcal}  [0,1]  \st 
    \forall t \in \Tcal. \, \sum_{(x,y,z) \in\Xcal\times\Ycal\times\Zcal} \theta_{x,y,z|t} =1 \rC \sm \Theta^{(0)}.
\end{align*}

We now split the data into two disjoint sets: $\Dbf = \Dbf^{(0)} \dcup \Dbf^{(1)}$. 
We use the notation $\Xbf^{(i)}$, $\Ybf^{(i)}$, $\Zbf^{(i)}$, $\Tbf^{(i)}$ for the corresponding components of $\Dbf^{(i)}$ for $i=0,1$.
Also let $I_i:=\lC n \in \lC 1,\dots,N \rC \st W_n:= (X_n,Y_n,Z_n,T_n) \in \Dbf^{(i)} \rC$ be the index set of $\Dbf^{(i)}$.
For brevity we put $\Wcal:=\Xcal\times\Ycal\times\Zcal\times\Tcal$.
We further make the following abbreviations for $(x,y,z,t) \in \Wcal$:
\begin{align*}
    N^{(i)}_{x,y,z,t} &:=  \sum_{(X_n,Y_n,Z_n,T_n) \in \Dbf^{(i)}} \I_{(x,y,z,t)}(X_n,Y_n,Z_n,T_n) \\ 
    &= \#\lC n \in I_i \st (X_n,Y_n,Z_n,T_n) = (x,y,z,t) \rC,
\end{align*}
and the following suggestive derived abbreviations:
\begin{align*}
    N^{(i)}_{x,y,z} &:=  \sum_{t \in \Tcal} N^{(i)}_{x,y,z,t}, &
    N^{(i)}_{x,z} &:=  \sum_{(y,t) \in \Ycal\times\Tcal} N^{(i)}_{x,y,z,t},\\
    N^{(i)}_{z} &:=  \sum_{(x,y,t) \in \Xcal\times\Ycal\times\Tcal} N^{(i)}_{x,y,z,t},&
    N^{(i)}_{t} &:=  \sum_{(x,y,z) \in \Xcal\times\Ycal\times\Zcal} N^{(i)}_{x,y,z,t},
\end{align*}
and so forth. It is then $N^{(i)} = \#\Dbf^{(i)}$.

We then write the mass functions as:
\[ p(x,y,z|t,\theta) = \prod_{(x',y',z',t')\in \Wcal} \lp \theta_{x',y',z'|t'} \rp^{\I_{(x',y',z',t')}(x,y,z,t)}.\]
The maximum likelihood objective for $H_1$ is then:
\begin{align*} \hat \theta^{(1)} := \hat \theta^{(1)}(\Dbf^{(1)}) 
&:= \argmax_{\theta \in\Theta^{(1)}} p(\Xbf^{(1)},\Ybf^{(1)},\Zbf^{(1)}|\Tbf^{(1)},\theta) \\
&= \argmax_{\theta \in\Theta^{(1)}} \prod_{n \in I_1}  p(X_n,Y_n,Z_n|T_n,\theta) \\
&= \argmax_{\theta \in\Theta^{(1)}} \prod_{n \in I_1} \prod_{(x,y,z,t)\in \Wcal} \lp \theta_{x,y,z|t} \rp^{\I_{(x,y,z,t)}(X_n,Y_n,Z_n,T_n)} \\
&= \argmax_{\theta \in\Theta^{(1)}} \prod_{(x,y,z,t)\in \Wcal} \prod_{n \in I_1} \lp \theta_{x,y,z|t} \rp^{\I_{(x,y,z,t)}(X_n,Y_n,Z_n,T_n)} \\
&= \argmax_{\theta \in\Theta^{(1)}} \prod_{(x,y,z,t)\in \Wcal} \lp \theta_{x,y,z|t} \rp^{\sum_{n \in I_1} \I_{(x,y,z,t)}(X_n,Y_n,Z_n,T_n)} \\
&= \argmax_{\theta \in\Theta^{(1)}} \prod_{(x,y,z,t)\in \Wcal} \lp \theta_{x,y,z|t} \rp^{ N_{x,y,z,t}^{(1)}}.% \\
%&= \argmax_{\theta \in\Theta^{(1)}} \prod_{\substack{(x,y,z,t)\in \Wcal\\N_{x,y,z,t}^{(1)}\neq 0}} \lp \theta_{x,y,z|t} \rp^{ N_{x,y,z,t}^{(1)}}\\
%&= \argmax_{\theta \in\Theta^{(1)}} \prod_{\substack{t \in \Tcal\\N_t^{(1)}\neq 0}} \prod_{\substack{(x,y,z,)\in \Xcal\times\Ycal\times\Zcal\\N_{x,y,z,t}^{(1)}\neq 0}} \lp \theta_{x,y,z|t} \rp^{ N_{x,y,z,t}^{(1)}}.
\end{align*}
This constrained optimization problem then leads to the following objective function with Lagrange multipliers $\lambda_t$, $t \in \Tcal$:
\[ \Lcal_1(\theta) := \sum_{(x,y,z,t)\in \Wcal} N_{x,y,z,t}^{(1)} \cdot \log \theta_{x,y,z|t} + 
\sum_{t \in \Tcal} \lambda_{t} \cdot \lp 1- \sum_{(x,y,z) \in\Xcal\times\Ycal\times\Zcal} \theta_{x,y,z|t}\rp.   \]
The KKT-conditions then demand for all $(x,y,z,t) \in \Wcal$:
\begin{align*}
    0= \frac{\partial \Lcal_1}{\partial\theta_{x,y,z|t}} &= \frac{N_{x,y,z,t}^{(1)}}{\theta_{x,y,z|t}}  - \lambda_t.
\end{align*}
If we sum for fixed $t \in \Tcal$ over all $(x,y,z) \in \Xcal\times\Ycal\times\Zcal$ we get:
\begin{align*}
    N_{t}^{(1)} &:= \sum_{(x,y,z) \in \Xcal\times\Ycal\times\Zcal} N_{x,y,z,t}^{(1)} 
        =  \lambda_t \cdot \sum_{(x,y,z) \in \Xcal\times\Ycal\times\Zcal} \theta_{x,y,z|t} 
        = \lambda_t \cdot 1
        =\lambda_t.
\end{align*}
Plugging $\lambda_t$ back into the equation above gives us the following frequency estimator for $\theta_{x,y,z|t}$ 
for every $(x,y,z,t) \in \Wcal$ with $N_{t}^{(1)} \neq 0$:
\[ \hat \theta_{x,y,z|t}^{(1)} = \frac{N_{x,y,z,t}^{(1)}}{N_{t}^{(1)}}.\]
For $N_{t}^{(1)} = 0$ there are no restrictions imposed by the maximum likelihood objective for setting $\theta_{x,y,z|t}$ besides normalization. 
E.g.\ we can use for the case $N_{t}^{(1)} = 0$:
\[ \hat \theta_{x,y,z|t}^{(1)} := \frac{N_{x,y,z}^{(1)}}{N^{(1)}}.\]
Note that these choices might effect the power of the split likelihood ratio estimator.

The maximum likelihood objective for $H_0$ is:
\begin{align*} \hat \theta^{(0)} &:= \hat \theta^{(0)}(\Dbf^{(0)}) \\
&:= \argmax_{\theta \in\Theta^{(0)}} p(\Xbf^{(0)},\Ybf^{(0)},\Zbf^{(0)}|\Tbf^{(0)},\theta) \\
&= \argmax_{\theta \in\Theta^{(0)}} \prod_{n \in I_0}  p(X_n,Y_n,Z_n|T_n,\theta) \\
&= \argmax_{\theta \in\Theta^{(0)}}  \prod_{n \in I_0}\prod_{(x,y,z,t)\in \Wcal} \lp \theta_{x|z} \cdot \theta_{y,z|t} \rp^{\I_{(x,y,z,t)}(X_n,Y_n,Z_n,T_n)} \\
&= \argmax_{\theta \in\Theta^{(0)}} \prod_{(x,y,z,t)\in \Wcal}\prod_{n \in I_0} \lp \theta_{x|z}  \rp^{\I_{(x,y,z,t)}(X_n,Y_n,Z_n,T_n)}\\ 
& \qquad \qquad \cdot \prod_{(x,y,z,t)\in \Wcal}\prod_{n \in I_0} \lp \theta_{y,z|t} \rp^{\I_{(x,y,z,t)}(X_n,Y_n,Z_n,T_n)} \\
&= \argmax_{\theta \in\Theta^{(0)}} \prod_{(x,y,z,t)\in \Wcal}\lp \theta_{x|z}  \rp^{N_{x,y,z,t}^{(0)}} \cdot \prod_{(x,y,z,t)\in \Wcal} \lp \theta_{y,z|t} \rp^{N_{x,y,z,t}^{(0)}} \\
&= \argmax_{\theta \in\Theta^{(0)}} \prod_{(x,z)\in \Xcal\times\Zcal}\lp \theta_{x|z}  \rp^{N_{x,z}^{(0)}} \cdot \prod_{(y,z,t)\in \Ycal\times\Zcal\times\Tcal} \lp \theta_{y,z|t} \rp^{N_{y,z,t}^{(0)}}.
\end{align*}
This constrained optimization problem then leads to the following objective function with Lagrange multipliers
$\lambda_z$ and $\lambda_t$ for $z \in \Zcal$ and $t \in \Tcal$, resp.:
\begin{align*}
    \Lcal_0(\theta) &:= \sum_{z\in \Zcal} \lB \sum_{x \in \Xcal}  N_{x,z}^{(0)} \cdot \log  \theta_{x|z}   
    + \lambda_z \cdot \lp 1- \sum_{x \in \Xcal} \theta_{x|z} \rp\rB \\ 
      &\quad   + \sum_{t\in \Tcal} \lB \sum_{(y,z)\in \Ycal\times\Zcal}   N_{y,z,t}^{(0)} \cdot \log \theta_{y,z|t} 
      + \lambda_t \cdot \lp 1 - \sum_{(y,z) \in \Ycal\times\Zcal} \theta_{y,z|t}  \rp \rB.
\end{align*}
The KKT-conditions then demand for all $(x,y,z,t) \in \Wcal$:
\begin{align*}
    0= \frac{\partial \Lcal_0}{\partial\theta_{x|z}} &= \frac{N_{x,z}^{(0)}}{\theta_{x|z}}  - \lambda_z, &
    0= \frac{\partial \Lcal_0}{\partial\theta_{y,z|t}} &= \frac{N_{y,z,t}^{(0)}}{\theta_{y,z|t}}  - \lambda_t.
\end{align*}
Again, summing up gives:
\begin{align*}
    N_z^{(0)} &= \sum_{x \in \Xcal} N_{x,z}^{(0)} = \lambda_z \cdot \sum_{x \in \Xcal} \theta_{x|z} = \lambda_z, &
    N_t^{(0)} &= \sum_{(y,z) \in\Ycal\times\Zcal} N_{y,z,t}^{(0)} = \lambda_t \cdot \sum_{x \in \Xcal} \theta_{y,z|t} = \lambda_t.
\end{align*}
Plugging the Lagrange multipliers back into the former equations gives us the frequency estimators
for $N_z^{(0)} \neq 0$ and $N_t^{(0)} \neq 0$, resp.:
\begin{align*}
    \hat \theta_{x|z}^{(0)} & = \frac{N_{x,z}^{(0)}}{N_z^{(0)}}, & \hat \theta_{y,z|t}^{(0)} & = \frac{N_{y,z,t}^{(0)}}{N_t^{(0)}}.
\end{align*}
If $N_z^{(0)} =0$ (if $N_t^{(0)} =0$, resp.), we can make the somewhat arbitrary choices:
\begin{align*}
    \hat \theta_{x|z}^{(0)} & := \frac{N_{x}^{(0)}}{N^{(0)}}, & \hat \theta_{y,z|t}^{(0)} & := \frac{N_{y,z}^{(0)}}{N^{(0)}}.
\end{align*}
With these derivations we can now come to our test statistic:
\begin{align*}
    E = E(\Dbf) 
      &= \frac{ p(\Xbf^{(0)},\Ybf^{(0)},\Zbf^{(0)}|\Tbf^{(0)},\hat\theta^{(1)}(\Dbf^{(1)})) }{p(\Xbf^{(0)},\Ybf^{(0)},\Zbf^{(0)}|\Tbf^{(0)},\hat\theta^{(0)}(\Dbf^{(0)}))}\\
   &= \prod_{(x,y,z,t) \in \Wcal} \lp  \frac{\hat\theta^{(1)}_{x,y,z,t} }{ \hat\theta^{(0)}_{x|z} \cdot \hat\theta^{(0)}_{y,z|t}} \rp^{N^{(0)}_{x,y,z,t}} \\
   &= \prod_{(x,y,z,t) \in \Wcal} \lp \frac{N_{x,y,z,t}^{(1)} \cdot N_z^{(0)} \cdot N_t^{(0)}}{N_{t}^{(1)} \cdot N_{x,z}^{(0)} \cdot N_{y,z,t}^{(0)} } \rp^{N^{(0)}_{x,y,z,t}},
\end{align*}
where we for simplicity only wrote down the case for $N_{t}^{(1)} \neq 0$. Note that in this estimator we do not need to resolve the ambiguity 
for the mentioned cases $N_z^{(0)} =0$ or $N_t^{(0)} =0$, as these do not occur in $E$, since then also $N^{(0)}_{x,y,z,t}=0$ and the corresponding term drops out.
With this we then get:
\begin{align*}
    \Lambda = \log E = \sum_{(x,y,z,t) \in \Xcal\times\Ycal\times\Zcal\times\Tcal}  N^{(0)}_{x,y,z,t} \cdot  \log \frac{N_{x,y,z,t}^{(1)} \cdot N_z^{(0)} \cdot N_t^{(0)}}{N_{t}^{(1)} \cdot N_{x,z}^{(0)} \cdot N_{y,z,t}^{(0)}}.
\end{align*}
\begin{comment}
If we want to include the ambiguous cases we could, more precisely, write:
\begin{align*}
    \Lambda 
    &= \sum_{\substack{(x,y,z,t) \in \Xcal\times\Ycal\times\Zcal\times\Tcal\\N^{(0)}_{x,y,z,t} \neq 0 \\N_{t}^{(1)} \neq 0}}  N^{(0)}_{x,y,z,t} \cdot  \log \frac{N_{x,y,z,t}^{(1)} \cdot N_z^{(0)} \cdot N_t^{(0)}}{N_{t}^{(1)} \cdot N_{x,z}^{(0)} \cdot N_{y,z,t}^{(0)}} \\
    & \qquad + \sum_{\substack{(x,y,z,t) \in \Xcal\times\Ycal\times\Zcal\times\Tcal\\N^{(0)}_{x,y,z,t} \neq 0 \\N_{t}^{(1)} = 0}}  N^{(0)}_{x,y,z,t} \cdot  \log \frac{N_{x,y,z}^{(1)} \cdot N_z^{(0)} \cdot N_t^{(0)}}{N^{(1)} \cdot N_{x,z}^{(0)} \cdot N_{y,z,t}^{(0)}}.
\end{align*}
\end{comment}
We then pick a significance level $\alpha \in [0,1]$ and declare $X$ dependent on $Y$ given $Z$ under $P(X,Y,Z|T)$ if:
\[ \Lambda \ge - \log \alpha.  \]
As discussed before, this rule leads to a \emph{non-asymptotic} type-I-error control of $\alpha$, 
which even holds for every \emph{joint} probability distribution on the input space $\Tcal^N$ that would have lead to the inputs.

\PF{

We should maybe assume that the marginal $P(Y,Z|T)$ is given for both, $H_0$ and $H_1$. Then we get the simpler estimator maximum-likelihood estimator $E$-variable:
\begin{align*}
    \Lambda = \log E = \sum_{(x,y,z,t) \in \Xcal\times\Ycal\times\Zcal\times\Tcal}  N^{(0)}_{x,y,z,t} \cdot  \log \frac{N_{x,y,z,t}^{(1)} \cdot N_z^{(0)}}{N_{y,z,t}^{(1)}\cdot N_{x,z}^{(0)}}.
\end{align*}

}

For finite discrete spaces we want to test:
\[ H_0:\, X \Indep_{P(X,Y,Z)} Y \given Z, \qquad \text{vs.} \qquad H_1:\,X \nIndep_{P(X,Y,Z)} Y \given Z. \]
Assume that the data comes sequentially in $K$ batches.
We apply the conditional independence test based on the $K$-split likelihood ratio test:
For each $k=1,\dots,K$ we have the logarithm of an E-variable:
\begin{align*}
     \log E_k := \sum_{(x,y,z) \in \Xcal\times\Ycal\times\Zcal}  N^{(k)}_{x,y,z} \cdot  \log \frac{N_{x,y,z}^{(<k)} \cdot N_z^{(k)} \cdot N^{(k)}}{N^{(<k)} \cdot N_{x,z}^{(k)} \cdot N_{y,z}^{(k)}},
\end{align*}
where $N_{x,y,z}^{(<k)}$ is just the number of counts of occurences of the tuple $(x,y,z)$ in the batches of data before $k$, etc., and $N_{x,y,z}^{(k)}$ the corresponding quantity counted just inside the batch with number $k$.
We can combine this into one logarithm of an E-variable:
\begin{align*}
    \Lambda := \log E :=  \sum_{k=1}^K \log E_k.
\end{align*}
For $\alpha \in [0,1]$ and declare $X$ dependent on $Y$ given $Z$ under $P(X,Y,Z)$ if:
\[ \Lambda \ge - \log \alpha.  \]


Conditional independence testing for finite discrete spaces using the $K$-split likelihood ratio test
 leads to the following log of an E-variable:
\begin{align*}
    \Lambda = \log E = \sum_{k=1}^K \sum_{(x,y,z,t) \in \Xcal\times\Ycal\times\Zcal\times\Tcal}  N^{(k)}_{x,y,z,t} \cdot  \log \frac{N_{x,y,z,t}^{(<k)} \cdot N_z^{(k)} \cdot N_t^{(k)}}{N_{t}^{(<k)} \cdot N_{x,z}^{(k)} \cdot N_{y,z,t}^{(k)}},
\end{align*}


\subsubsection{Linear Gaussian Conditional Random Variables}

We can similarly apply the split likelihood ratio test for conditional independence testing 
to linear Gaussian conditional random variables, 
or other conditional random variables with tractable parametric Markov kernels.

\begin{Lem}
    \label{lem:mle-lin-gauss}
    Consider a linear Gaussian Markov kernel $P(X|Z)$ with density:
    \begin{align*}
        p(x|z) &=\Ncal\lp x \st \Gamma \cdot z + \gamma, \Sigma \rp.
    \end{align*}
    Further assume that we have i.i.d.\ data 
    $\Dbf=\lC (X_1,Z_1),\dots,(X_N,Z_N) \rC$ from: 
    \[ P^{\otimes N}(\Xbf|\Zbf) := \bigotimes_{n=1}^N P(X_n|Z_n),  \]
    where $P(X_n|Z_n)$ is a copy of $P(X|Z)$.
    Then the maximum likelihood estimators for $\Gamma$, $\gamma$, $\Sigma$ are given by:
    \begin{align*}
        \hat \Gamma &:= \hat \Sigma_{X,Z} \hat \Sigma_{Z,Z}^{-1}, &
        \hat \gamma &:= \hat \mu_X -  \hat \Sigma_{X,Z} \hat \Sigma_{Z,Z}^{-1}\, \hat \mu_Z, &
        \hat \Sigma &:= \hat\Sigma_{X,X} -\hat\Sigma_{X,Z}\,\hat\Sigma_{Z,Z}^{-1}\,\hat\Sigma_{Z,X},
    \end{align*}
  where:
      \begin{align*}
       \hat \mu_X &:= \frac{1}{N} \sum_{n=1}^N X_n, &
       \hat \mu_Z &:= \frac{1}{N} \sum_{n=1}^N Z_n, \\
       \hat \Sigma_{X,X} &:= \lp \frac{1}{N} \sum_{n=1}^N X_n \cdot X_n^\top \rp - \hat \mu_X \cdot \hat \mu_X^\top, &
        \hat \Sigma_{Z,Z} &:= \lp \frac{1}{N} \sum_{n=1}^N Z_n \cdot Z_n^\top \rp - \hat \mu_Z \cdot \hat \mu_Z^\top,\\
        \hat \Sigma_{X,Z} &:= \lp \frac{1}{N} \sum_{n=1}^N X_n \cdot Z_n^\top \rp - \hat \mu_X \cdot \hat \mu_Z^\top, &
        \hat \Sigma_{Z,X} &:= \hat \Sigma_{X,Z}^\top.
    \end{align*}
    In other words, this maximum likelihood estimator agrees with the one for a \emph{jointly} Gaussian distributed $(X,Z)$
    and then using the conditioning formula for Gaussians conditioned on $Z$. %In particular, it implicitely treats $Z$ as a Gaussian distributed random variable.
    \begin{proof}
        Let $\theta=(\Gamma,\gamma,\Sigma)$. The maximum likelihood estimate for $\theta$ is then given by:
    \begin{align*}
        \hat \theta & =\argmax_{\theta} \prod_{n=1}^N \Ncal\lp X_n \st \Gamma \cdot Z_n + \gamma, \Sigma \rp \\
        &= \argmin_{\theta} - 2 \log  \prod_{n=1}^N \Ncal\lp X_n \st \Gamma \cdot Z_n + \gamma, \Sigma \rp \\
        &= \argmin_{\theta} \lB N \cdot \log |\det(2\pi\Sigma)| + \sum_{n=1}^N \Tr\lp \Sigma^{-1}(\Gamma \cdot Z_n + \gamma - X_n)(\Gamma \cdot Z_n + \gamma - X_n)^\top \rp \rB.
    \end{align*}
Let the latter expression be abbreviate with:
\[\Lcal  :=N \cdot \log |\det(2\pi\Sigma)| + \sum_{n=1}^N \Tr\lp \Sigma^{-1}(\Gamma \cdot Z_n + \gamma - X_n)(\Gamma \cdot Z_n + \gamma - X_n)^\top \rp.\]
    We will then take derivatives and equal them to zero to get the wanted expressions. 
Note that we will use the following general matrix identities, see \cite{matrixcookbook}, with $\Dbf$ symmetric:
\begin{align*}
    %\frac{\partial \det(\Abf \Xbf \Bbf)}{\partial \Xbf} &= \det(\Abf \Xbf \Bbf) \cdot \Xbf^{-\top},\\
    \frac{\partial \log | \det(\Xbf)|}{\partial \Xbf} &= \Xbf^{-\top},\\
    \frac{\partial \Tr(\Abf \Xbf^{-1} \Bbf) }{ \partial \Xbf } &= - \lp \Xbf^{-1} \Bbf \Abf \Xbf^{-1} \rp^\top,\\
    %\frac{\partial \Tr(\Bbf \Xbf \Xbf^\top) }{ \partial \Xbf } &= \Bbf\Xbf + \Bbf^\top \Xbf \\
    %\frac{\partial}{\partial \Xbf} \Tr \lp  (\Abf\Xbf\Bbf+\Cbf)^\top \Dbf (\Abf\Xbf\Bbf+\Cbf)\rp 
    \frac{\partial}{\partial \Xbf} \Tr \lp \Dbf (\Abf\Xbf\Bbf+\Cbf)(\Abf\Xbf\Bbf+\Cbf)^\top\rp 
    &= 2 \Abf^\top \Dbf (\Abf \Xbf \Bbf + \Cbf) \Bbf^\top.
\end{align*}
We get for $\gamma$:
\begin{align*}
 0 = \frac{\partial \Lcal}{\partial \gamma} & = \sum_{n=1}^N  2 \cdot \Sigma^{-1}\,(\Gamma \cdot Z_n + \gamma - X_n) \\
 &= 2 \cdot N \cdot \Sigma^{-1} \, ( \Gamma\, \hat \mu_Z + \gamma - \hat \mu_X), 
 \end{align*}
 which implies:
 \[ \Gamma\, \hat \mu_Z + \gamma - \hat \mu_X = 0.  \]
We get for $\Gamma$:
\begin{align*}
 0 = \frac{\partial \Lcal}{\partial \Gamma} & = \sum_{n=1}^N  2 \cdot \Sigma^{-1}\,(\Gamma \cdot Z_n + \gamma - X_n)\, Z_n^\top \\
 &= 2 \cdot \Sigma^{-1} \lp \Gamma \sum_{n=1}^N Z_n Z_n^\top + \gamma \cdot \sum_{n=1}^N Z_n - \sum_{n=1}^N X_n Z_n^\top    \rp\\ 
 &= 2 \cdot N \cdot \Sigma^{-1} \lp \Gamma\, \hat \Sigma_{Z,Z} + \Gamma \hat\mu_Z \hat \mu_Z^\top + \gamma \cdot \hat \mu_Z - 
  \hat \Sigma_{X,Z} - \hat \mu_X \hat \mu_Z^\top \rp, 
 \end{align*}
 which implies:
 \begin{align*}
     0 &= \Gamma\, \hat \Sigma_{Z,Z} - \hat \Sigma_{X,Z} + \lp \Gamma \hat\mu_Z  + \gamma - \hat \mu_X \rp \hat \mu_Z^\top \\
       &= \Gamma\, \hat \Sigma_{Z,Z} - \hat \Sigma_{X,Z}.
 \end{align*}
Solving for $\Gamma$ and then $\gamma$ gives the maximum likelihood estimators:
\begin{align*}
    \hat\Gamma & = \hat\Sigma_{X,Z} \,\hat\Sigma_{Z,Z}^{-1}, \\
    \hat\gamma & = \hat \mu_X - \hat\Gamma\,\hat\mu_Z = \hat \mu_X - \hat\Sigma_{X,Z} \,\hat\Sigma_{Z,Z}^{-1}\,\hat\mu_Z.
\end{align*}
Now lets look at $\Sigma$:
\begin{align*}
    0 = \frac{\partial \Lcal}{\partial\Sigma} & = N \cdot \Sigma^{-1} - \sum_{n=1}^N  \Sigma^{-1} (\Gamma \cdot Z_n + \gamma - X_n)(\Gamma \cdot Z_n + \gamma - X_n)^\top \Sigma^{-1},
 \end{align*}
which leads to:
\begin{align*}
    \Sigma &= \frac{1}{N} \sum_{n=1}^N  (\Gamma \cdot Z_n + \gamma - X_n)(\Gamma \cdot Z_n + \gamma - X_n)^\top \\
           &= \frac{1}{N} \sum_{n=1}^N  \Gamma \cdot Z_n \cdot Z_n^\top \cdot \Gamma^\top + \Gamma \cdot Z_n \cdot \gamma^\top - 
           \Gamma \cdot Z_n \cdot X_n^\top \\ 
          &\qquad + \gamma \cdot X_n^\top \cdot \Gamma^\top + \gamma \cdot \gamma^\top - \gamma \cdot X_n^\top   
           - X_n \cdot Z_n^\top \cdot \Gamma^\top - X_n \cdot \gamma^\top + X_n \cdot X_n^\top  \\
          &= \Gamma \cdot ( \hat\Sigma_{Z,Z} + \hat\mu_Z \hat\mu_Z^\top) \cdot \Gamma^\top + \Gamma \cdot \hat\mu_Z\cdot\gamma^\top
          - \Gamma \cdot \hat\Sigma_{Z,X} - \Gamma \cdot \hat\mu_Z\hat\mu_X^\top \\
          &\qquad + \gamma \cdot \hat\mu_X^\top \cdot \Gamma^\top + \gamma \gamma^\top - \gamma \cdot \hat\mu_X^\top
          - \hat\Sigma_{X,Z} \cdot \Gamma^\top - \hat\mu_X\hat\mu_Z^\top \cdot \Gamma^\top 
          - \hat\mu_X \cdot \gamma^\top + \hat\Sigma_{X,X} + \hat\mu_X\hat\mu_X^\top.
 \end{align*}
If we now plug in our estimators from before we get for the expressions:
\begin{align*}
 \hat\Gamma\cdot\hat\Sigma_{Z,Z}\cdot\hat\Gamma^\top = \hat\Sigma_{X,Z}\cdot\hat\Gamma^\top = \hat\Gamma\cdot\hat\Sigma_{Z,X}
 = \hat\Sigma_{X,Z}\,\hat\Sigma_{Z,Z}^{-1}\,\hat\Sigma_{Z,X}. 
\end{align*}
We abbeviate in the following: $\hat \nu_X := \hat\Gamma\,\hat\mu_Z$. So $\hat \gamma = \hat\mu_X - \hat\nu_X$.
With this we get:
\begin{align*}
    \hat \Sigma &= \lp \hat\Sigma_{X,X} -\hat\Sigma_{X,Z}\,\hat\Sigma_{Z,Z}^{-1}\,\hat\Sigma_{Z,X} \rp 
     + \hat\nu_X\hat\nu_X^\top + \hat\nu_X (\hat\mu_X-\hat\nu_X) - \hat\nu_X\hat\mu_X^\top \\ 
    &\qquad + (\hat\mu_X-\hat\nu_X) \hat\nu_X^\top + (\hat\mu_X-\hat\nu_X)(\hat\mu_X-\hat\nu_X)^\top - 
    (\hat\mu_X-\hat\nu_X) \mu_X^\top  \\
    &\qquad - \hat\mu_X \hat\nu_X^\top - \hat\mu_X (\hat\mu_X-\hat\nu_X)^\top  + \hat\mu_X\hat\mu_X^\top\\
    & =\hat\Sigma_{X,X} -\hat\Sigma_{X,Z}\,\hat\Sigma_{Z,Z}^{-1}\,\hat\Sigma_{Z,X}.
\end{align*}
So our maximum likelihood estimators are given by:
\begin{align*}
    \hat\Gamma & = \hat\Sigma_{X,Z} \,\hat\Sigma_{Z,Z}^{-1}, \\
    \hat\gamma & = \hat \mu_X - \hat\Sigma_{X,Z} \,\hat\Sigma_{Z,Z}^{-1}\,\hat\mu_Z, \\
    \hat \Sigma &= \hat\Sigma_{X,X} -\hat\Sigma_{X,Z}\,\hat\Sigma_{Z,Z}^{-1}\,\hat\Sigma_{Z,X}.
\end{align*}
This shows the claim.
    \end{proof}
\end{Lem}


\begin{Rem}[Split likelihood ratio test for conditional independence of linear Gaussian conditional random variables]
  Consider three conditional random variables $X$, $Y$, $Z$ that come from an underlying linear Gaussian Markov kernel $P(X,Y,Z|T)$.
  We want to test the following conditional independence:
  \[H_0:\; X \Indep_{P(X,Y,Z|T)} Y \given Z, \qquad \text{vs.} \qquad H_1:\; X \nIndep_{P(X,Y,Z|T)} Y \given Z, \]
  based on i.i.d.\ data $\Dbf = \lC (X_1,Y_1,Z_1,T_1),\dots,(X_N,Y_N,Z_N,T_N) \rC$ from $P(X,Y,Z|T)$.
  We split the data set into two disjoint subsets $\Dbf = \Dbf^{(0)} \dcup \Dbf^{(1)}$.
  We then consider the following test statistic, following the split likelihood ratio test:
 % \[ \Lambda := \log \frac{\Ncal\lp\Xbf^{(0)},\Ybf^{(0)},\Zbf^{(0)}\st \hat\Gamma^{(1)}\cdot \Tbf^{(0)}+\hat\gamma^{(1)},\hat\Sigma^{(1)}\rp}{ \Ncal\lp\Xbf^{(0)}\st \hat\Gamma_{X|Z}^{(0)}\cdot \Zbf^{(0)}+\hat\gamma_{X|Z}^{(0)},\hat\Sigma_{X|Z}^{(0)}\rp \cdot \Ncal\lp\Ybf^{(0)},\Zbf^{(0)}\st \hat\Gamma^{(0)}\cdot \Tbf^{(0)}+\hat\gamma^{(0)},\hat\Sigma^{(0)}\rp },  \]
  \[ \Lambda := \sum_{n \in I_0}  %\sum_{(X_n,Y_n,Z_n,T_n) \in \Dbf^{(0)}}
  \log \frac{\Ncal\lp \matb{X_n&Y_n&Z_n}^\top\st \hat\Gamma^{(1)}\cdot T_n+\hat\gamma^{(1)},\hat\Sigma^{(1)}\rp}{ \Ncal\lp X_n\st \hat\Gamma_{X|Z}^{(0)}\cdot Z_n+\hat\gamma_{X|Z}^{(0)},\hat\Sigma_{X|Z}^{(0)}\rp \cdot \Ncal\lp \matb{Y_n&Z_n}^\top\st \hat\Gamma^{(0)}\cdot T_n+\hat\gamma^{(0)},\hat\Sigma^{(0)}\rp },  \]
  with the corresponding maximum likelihood estimators based on the empirical means $\hat\mu^{(i)}$ and covariance matrices $\hat\Sigma^{(i)}$, see Lemma \ref{lem:mle-lin-gauss}:
  \begin{align*}
  \hat\Gamma_{X|Z}^{(0)} &:= \hat\Sigma_{X,Z}^{(0)} \,{\hat\Sigma_{Z,Z}^{(0),-1}},\\
  \hat\gamma_{X|Z}^{(0)} &:= \hat\mu_X^{(0)} - \hat\Sigma_{X,Z}^{(0)} \,{\hat\Sigma_{Z,Z}^{(0),-1}}\,\hat\mu_Z^{(0)},\\
  \hat\Sigma_{X|Z}^{(0)} &:= \hat\Sigma_{X,X}^{(0)} -\hat\Sigma_{X,Z}^{(0)}\,{\hat\Sigma_{Z,Z}^{(0),-1}}\,\hat\Sigma_{Z,X}^{(0)},\\
      \hat\Gamma^{(0)} &:= \matb{\hat\Sigma_{Y,T}^{(0)}\\\hat\Sigma_{Z,T}^{(0)}} \,{\hat\Sigma_{T,T}^{(0),-1}},\\
      \hat\gamma^{(0)} &:= \matb{\hat\mu_Y^{(0)}\\\hat\mu_Z^{(0)}} - \matb{\hat\Sigma_{Y,T}^{(0)}\\\hat\Sigma_{Z,T}^{(0)}} \,{\hat\Sigma_{T,T}^{(0),-1}}\,\hat\mu_T^{(0)}, \\
     \hat\Sigma^{(0)} &:= \matb{\hat\Sigma_{Y,Y}^{(0)}&\hat\Sigma_{Y,Z}^{(0)}\\\hat\Sigma_{Z,Y}^{(0)}&\hat\Sigma_{Z,Z}^{(0)}} - \matb{\hat\Sigma_{Y,T}^{(0)}\\\hat\Sigma_{Z,T}^{(0)}}\,
      {\hat\Sigma_{T,T}^{(0),-1}}\,\matb{\hat\Sigma_{T,Y}^{(0)}&\hat\Sigma_{T,Z}^{(0)}},\\
      \hat\Gamma^{(1)} &:= \matb{\hat\Sigma_{X,T}^{(1)}\\\hat\Sigma_{Y,T}^{(1)}\\\hat\Sigma_{Z,T}^{(1)}} \,{\hat\Sigma_{T,T}^{(1),-1}},\\
      \hat\gamma^{(1)} &:= \matb{\hat\mu_X^{(1)}\\\hat\mu_Y^{(1)}\\\hat\mu_Z^{(1)}} - \matb{\hat\Sigma_{X,T}^{(1)}\\\hat\Sigma_{Y,T}^{(1)}\\\hat\Sigma_{Z,T}^{(1)}} \,{\hat\Sigma_{T,T}^{(1),-1}}\,\hat\mu_T^{(1)}, \\
     \hat\Sigma^{(1)} &:= \matb{\hat\Sigma_{X,X}^{(1)}&\hat\Sigma_{X,Y}^{(1)}&\hat\Sigma_{X,Z}^{(1)}\\\hat\Sigma_{Y,X}^{(1)}&\hat\Sigma_{Y,Y}^{(1)}&\hat\Sigma_{Y,Z}^{(1)}\\\hat\Sigma_{X,Z}^{(1)}&\hat\Sigma_{Z,Y}^{(1)}&\hat\Sigma_{Z,Z}^{(1)}} 
     - \matb{\hat\Sigma_{X,T}^{(1)}\\\hat\Sigma_{Y,T}^{(1)}\\\hat\Sigma_{Z,T}^{(1)}} \,{\hat\Sigma_{T,T}^{(1),-1}}\,
     \matb{\hat\Sigma_{T,X}^{(1)}&\hat\Sigma_{T,Y}^{(1)}&\hat\Sigma_{T,Z}^{(1)}}.
  \end{align*}
  We then declare $X$ dependent on $Y$ given $Z$ w.r.t.\ $P(X,Y,Z|T)$ if:
  \[ \Lambda \ge -\log \alpha, \]
  for some pre-specified significance level $\alpha \in [0,1]$, which bounds the type-I-error of this test.
\end{Rem}


\PF{
\subsubsection{Some Thoughts}

Why do we, in causal inference/graph learning test:
  \[H_0:\; X \Indep_{P(X,Y,Z|T)} Y \given Z, \qquad \text{vs.} \qquad H_1:\; X \nIndep_{P(X,Y,Z|T)} Y \given Z, \]
  instead of the reverse order:
 \[H_0:\; X \nIndep_{P(X,Y,Z|T)} Y \given Z, \qquad \text{vs.} \qquad H_1:\; X \Indep_{P(X,Y,Z|T)} Y \given Z. \]
???

For sake of arguments let $Z=\ast$.

I do understand that in statistics one can consider $P(X) \otimes P(Y)$ a high-entropy or noisy version of the joint $P(X,Y)$.
So if one wants to find a signal in the data as a statistical relationship between $X$ and $Y$ then one would need to provide strong evidence that the signal is actually there, so testing:
  \[H_0:\; X \Indep Y, \qquad \text{vs.} \qquad H_1:\; X \nIndep Y, \]
makes sense to me in this context.

However, when we want to learn a graph then by default I would assume that all variables are somehow related 
and (conditional) independencies would remove edges. So in this context I would rather expect that one would need 
to present strong evidence that one can safely remove edges, via (conditional) independence tests. So here I would rather by default
test the reverse order:
  \[H_0:\; X \nIndep Y, \qquad \text{vs.} \qquad H_1:\; X \Indep Y. \]
Does this make sense to you?

I do understand that getting the asymptotic thresholds right is easier under independence, i.e.\ if $H_0:\; X \Indep Y$, but I would think that this is not justification enough.

The split likelihood ratio test at least does not require an asymptotic threshold and at least here one could just test the reverse:
 \[H_0:\; X \nIndep_{P(X,Y,Z|T)} Y \given Z, \qquad \text{vs.} \qquad H_1:\; X \Indep_{P(X,Y,Z|T)} Y \given Z. \]
In the categorical setting we would just use:
\begin{align*}
    \Lambda := \sum_{(x,y,z,t) \in \Xcal\times\Ycal\times\Zcal\times\Tcal}  N^{(0)}_{x,y,z,t} \cdot  \log \frac{N_{t}^{(0)} \cdot N_{x,z}^{(1)} \cdot N_{y,z,t}^{(1)}}{N_{x,y,z,t}^{(0)} \cdot N_z^{(1)} \cdot N_t^{(1)}},
\end{align*}
and declare $X$ independent from $Y$ given $Z$ w.r.t.\ $P(X,Y,Z|T)$ if:
\[ \Lambda \ge - \log \alpha. \]
With this test the error that we would confuse weak dependence with an independence is bounded by $\alpha$.

Similarly, for linear Gaussians variable we could do the same by using:
  \[ \Lambda := \sum_{n \in I_0}    \log \frac{ \Ncal\lp X_n\st \hat\Gamma_{X|Z}^{(1)}\cdot Z_n+\hat\gamma_{X|Z}^{(1)},\hat\Sigma_{X|Z}^{(1)}\rp \cdot \Ncal\lp \matb{Y_n&Z_n}^\top\st \hat\Gamma^{(1)}\cdot T_n+\hat\gamma^{(1)},\hat\Sigma^{(1)}\rp }{\Ncal\lp \matb{X_n&Y_n&Z_n}^\top\st \hat\Gamma^{(0)}\cdot T_n+\hat\gamma^{(0)},\hat\Sigma^{(0)}\rp}. \]

Any thoughts??
}

\subsubsection{Type-I and Type-II Error for Testing based on E-Variables}


\begin{Thm}[Sanov's Theorem, see \cite{Csi84,Bal2020}]
    \label{sanov}
    Let $(\Xcal,P)$ be a probability space,  $X_1,\dots,X_N \sim P$ i.i.d.\ random variables and $P_N:=\frac{1}{N} \sum_{n=1}^N \delta_{X_n}$ the empirical distribution. 
    For any measurable set $\Acal \ins \Pcal(\Xcal)$ of probability measures on $\Xcal$ we denote by $\Acal^\cx=\{\E_{\mu \sim Q}[\mu]\,|\,Q \in \Pcal(\Acal) \}$ its complete convex hull.
    Then we have the inequality:
    \[ P\left( P_N \in \Acal \right) \le \exp\left( -N \cdot \inf_{Q \in \Acal^\cx} \KL(Q\|P) \right).   \]
\end{Thm}


\begin{Not}
    Let $E:\, \Xcal \to \R_{\ge0}$ be a measurable map. We then put for $\gamma, \delta \in \R_{>0}$:
    \begin{align*}
        \Hcal_0^E & := \lC Q \in \Pcal(\Xcal) \st \E_Q[E] \le 1 \rC,\\
        \Acal_0^E & := \lC Q \in \Pcal(\Xcal) \st \E_Q[\log E] \le 0 \rC,\\
        \Acal_\gamma^E & := \lC Q \in \Pcal(\Xcal) \st \E_Q[\log E] \le \gamma \rC,\\
        \Hcal_{1}^E(\gamma,\delta) & := \lC P \in \Pcal(\Xcal) \st \KL(\Acal_\gamma^E \|P) \ge \delta \rC,\\
        \Hcal^E_1& := \bigcup_{\gamma,\delta >0} \Hcal_{1}^E(\gamma,\delta).
    \end{align*}
    Then we have:
    \[  \Hcal_0^E \ins \Acal_0^E \ins \Acal_\gamma^E, \qquad \Hcal_{1}^E(\gamma,\delta) \ins \Hcal_1^E \ins \lp \Acal_0^E \rp^\cmpl.  \]
    Also note that for $0<\gamma_2 \le \gamma_1$ and $0<\delta_2 \le \delta_1$ we have the inclusions:
    \[ \Acal_{\gamma_2}^E  \ins \Acal_{\gamma_1}^E, \qquad \Hcal_{1}^E(\gamma_2,\delta_2)  \sni \Hcal_{1}^E(\gamma_1,\delta_1).   \]
\end{Not}


\begin{Thm}[Type-I and type-II error convergence]
    \label{thm:e-variables-type-I-II-error}
    %Let $\Hcal_0 \ins \Pcal(\Xcal)$ be a hypothesis class and $E$ an E-variable for $\Hcal_0$.
    Let $E:\, \Xcal \to \R_{\ge0}$ be a measurable map and for $n, N \in \N$ let $X_n:\, \Xcal^N \to \Xcal$ be the canonical projection map onto the $n$-th factor.
    % Let $(X_n)_{n \in \N} \sim P_*$ be an i.i.d.\ sequence of samples from an underlying true distribution $P_* \in \Pcal(\Xcal)$.
     For $N \in \N$ we abbreviate: % the product of E-variables evaluated at $X_n$ for $n=1,\dots,N$ as:
     \begin{align*} E_N:=E(X_{1:N}):=\prod_{n =1}^N E(X_n):\; \Xcal^N \to \R_{\ge 0}. \end{align*}
 Let $(\alpha_n)_{n \in \N} \ins \R_{>0}$ be sequence of positive real numbers.
 \begin{enumerate}
     \item  If the sequence converges to zero:
 \begin{align*} 
     \alpha_N &\stackrel{N \to \infty}{\longrightarrow} 0,
 \end{align*}
then we have the following uniform \emph{type-I-error} control:
 \begin{align*}
     \sup_{P \in \Hcal_0^E}  P^{\otimes N}\lp\log E_N \ge -\log\alpha_N \rp & \le \alpha_N \stackrel{N \to \infty}{\longrightarrow} 0.
 \end{align*}
\item If, instead or in addition, the sequence satisfies the following convergence condition:
 \begin{align*} 
     \gamma_N:=-\frac{1}{N} \log \alpha_N &\stackrel{N \to \infty}{\longrightarrow} 0, 
 \end{align*}
then we get the following:
\begin{enumerate}
 \item For every $P \in \Hcal_1^E$, i.e.\ with $\delta_P:=\KL(\Acal_\gamma^E\|P) >0$ for some $\gamma >0$,
 we have the following \emph{type-II-error} control:
 \begin{align*}
     P^{\otimes N}\lp \log E_N \le-\log\alpha_N \rp & \stackrel{N \ge N_\gamma}{\le} \exp\lp -N \cdot \delta_P \rp \stackrel{N \to \infty}{\longrightarrow} 0.
 \end{align*}
 \item If $\Hcal_1 \ins \Pcal(\Xcal)$ is an (alternative) hypothesis class inside $\Hcal_{1}^E(\gamma,\delta)$ for some $\gamma,\delta >0$, i.e.\ we have:
    \begin{align*}
         \inf_{P \in \Hcal_1} \inf_{Q \in \Acal_{\gamma}^E} \KL(Q \|P) \ge \delta > 0,
    \end{align*}
    then we have the following uniform \emph{type-II-error} control:
 \begin{align*}
     \sup_{P \in \Hcal_1}  P^{\otimes N}\lp \log E_N \le-\log\alpha_N \rp &   \stackrel{N \ge N_\gamma}{\le} \exp\lp -N \cdot \delta \rp \stackrel{N \to \infty}{\longrightarrow} 0.
 \end{align*}
\end{enumerate}
  Note that the indicated inequalities only hold for $N$ above some threshold $N_\gamma$.
  \begin{enumerate}[resume]
      \item If we, furthermore, have a sequence $(\delta_n)_{n \in \N} \ins \R_{>0}$ such that:
         \begin{align*} 
                 \delta_N &\stackrel{N \to \infty}{\longrightarrow} 0, &N \cdot \delta_N &\stackrel{N \to \infty}{\longrightarrow} \infty, 
         \end{align*}
       then we have the following uniform \emph{type-II-error} control with increasing alternative hypothesis class:
        \begin{align*}
            \sup_{P \in \Hcal_{1}^E(\gamma_N,\delta_N)}  P^{\otimes N}\lp \log E_N \le-\log\alpha_N \rp & \le \exp\lp -N \cdot \delta_N \rp \stackrel{N \to \infty}{\longrightarrow} 0.
        \end{align*}
    Note that the inequality here holds for every $N \in \N$ and that, due to the convergence
    $\gamma_N,\delta_N \stackrel{N \to \infty}{\longrightarrow} 0$, we have: 
    \[ \Hcal^E_1 = \bigcup_{N \in \N} \Hcal_{1}^E(\gamma_N,\delta_N).\] 
 \end{enumerate}
\end{enumerate}
\begin{proof}
    1.) The product of independent E-variables is again an E-variable, indeed for $P \in \Hcal_0^E$:
     \begin{align*}
         \E_P\lB E_N \rB =  \E_P\lB \prod_{n=1}^N E(X_n) \rB \stackrel{\Indep}{=} \prod_{n=1}^N \E_P\lB E(X_n) \rB \le \prod_{n=1}^N 1 = 1. 
    \end{align*}
    So for every $\alpha \in [0,1]$ we get by the Markov inequality for $P \in \Hcal_0^E$:
    \begin{align*}
        P\lp\log E_N \ge -\log\alpha \rp & = P\lp  E_N \ge \alpha^{-1} \rp \le \frac{\E_P\lB E_N \rB}{\alpha^{-1}} \le \alpha.
    \end{align*}
    If we now choose $\alpha=\alpha_N$ and use our assumption $\alpha_N \to 0$ then the type-I-error convergence to $0$ is clear.

    2.a) The case $P \in \Hcal_1^E$ follows from the next one with $\Hcal_1=\lC P \rC$.

    2.b) So consider a hypothesis class $\Hcal_1 \ins \Pcal(\Xcal)$ such that:
    \begin{align*}
       \inf_{P \in \Hcal_1} \KL(\Acal_\gamma^E \|P) \ge \delta  > 0,
    \end{align*}
    for some  $\gamma,\delta > 0$.
    Since by assumption $\gamma_N:=-\frac{1}{N} \log \alpha_N \to 0$ there exists an $N_\gamma \in \N$ such that for all $N \ge N_\gamma$ we have:
    \begin{align*}
      0 < \gamma_N <  \gamma.
    \end{align*}
    This shows that $\Acal_{\gamma_N} \ins \Acal_{\gamma}$  for all $N \ge N_\gamma$ and thus for all $P \in \Hcal_1$:
    \begin{align*}
        \KL(\Acal_{\gamma_N}^E \|P) \ge  \KL(\Acal_\gamma^E \|P) \ge \delta > 0.
    \end{align*}
    With this and Sanov's Theorem \ref{sanov}, see \cite{Csi84,Bal2020}, we get with $P_N := \frac{1}{N} \sum_{n=1}^N \delta_{X_n}$:
    \begin{align*}
        P\lp  \log E_N \le - \log(\alpha_N)  \rp 
      &=  P\lp   \sum_{n=1}^N  \log E(X_n)  \le - \log(\alpha_N)  \rp  \\
      &=  P\lp   \frac{1}{N}\sum_{n=1}^N  \log E(X_n)  \le - \frac{1}{N} \log(\alpha_N)  \rp  \\
      &=  P\lp  \E_{P_N}\lB \log(E) \rB \le \gamma_N  \rp  \\
      &=  P\lp P_N \in \Acal^E_{\gamma_N}\rp \\ 
      &\le \exp\lp -N \cdot \KL(\Acal^E_{\gamma_N}\|P) \rp \\
      &\stackrel{N\ge N_\gamma}{\le} \exp\lp -N \cdot \delta \rp\\
      & \longrightarrow 0,
    \end{align*}
    where the last inequality holds for all $N \ge N_\gamma$.
    Since the inequality holds for every $P \in \Hcal_1$ we get:
    \begin{align*}
        \sup_{P \in \Hcal_1}   P\lp  \log E_N \le - \log(\alpha_N)  \rp 
      &\stackrel{N\ge N_\gamma}{\le} \exp\lp -N \cdot \delta \rp  \longrightarrow 0.
    \end{align*}
    This shows the claims.

    2.c) The same arguments as above apply for $\delta=\delta_N$, $\gamma=\gamma_N$, $N=N_{\gamma_N}$. Note that we get the inequalities for every fixed $N$, not just in the limit.
 \end{proof}
\end{Thm}

\begin{Rem}
    The convergence property of the sequence $(\alpha_n)_{n \in \N} \ins \R_{>0}$ in Theorem \ref{thm:e-variables-type-I-II-error} trades off between
    type-I and type-II error control. We have all three cases:
    \begin{enumerate}
        \item The sequence $\alpha_N:=\frac{\alpha}{N^r}$ for fixed $\alpha,r >0$ satifies both conditions:
             \begin{align*} 
                 \alpha_N &\stackrel{N \to \infty}{\longrightarrow} 0, &\gamma_N:=-\frac{1}{N} \log \alpha_N &\stackrel{N \to \infty}{\longrightarrow} 0. 
         \end{align*}
     \item The sequence $\alpha_N:=\frac{1}{2^N}$ satifies the first, but not the second condition.
     \item The constant sequence $\alpha_N:=\alpha$ for fixed $\alpha > 0$ satisfies the second, but not the first condition.
    \end{enumerate}
%    For the sequence $(\delta_n)_{n \in \N} \ins \R_{>0}$ we can make the following choice:
%    \[ \delta_N := \frac{1}{\sqrt{N}}.  \]
\end{Rem}

Based on the result of Theorem \ref{thm:e-variables-type-I-II-error} we make the following definition:

\begin{Def}[Effective E-variable]
 Let $\Hcal_0,\Hcal_1 \ins \Pcal(\Xcal)$ be two hypothesis classes and $E:\, \Xcal \to \R_{\ge0}$ be a measurable map.
 We call $E$ an \emph{effective E-variable} for $(\Hcal_0,\Hcal_1)$ if:
 \[ \Hcal_0 \ins \Hcal_0^E \qquad \text{ and } \qquad \Hcal_1 \ins \Hcal_1^E. \]
\end{Def}

